\documentclass[11pt]{article}

\usepackage[margin=1in]{geometry}
\usepackage{amsmath,amssymb,amsthm,bm,mathtools}
\usepackage{booktabs}
\usepackage{array}
\usepackage{longtable}
\usepackage{tabularx}
\usepackage{hyperref}

\hypersetup{
  colorlinks=true,
  linkcolor=blue,
  urlcolor=blue,
  citecolor=blue
}

\setlength{\parindent}{0pt}
\setlength{\parskip}{0.7em}

\newcommand{\pf}{\mathbf{pf}}
\newcommand{\Xm}{\mathbf{X}_m}
\newcommand{\ym}{\mathbf{y}_m}
\newcommand{\zm}{\mathbf{z}_m}
\newcommand{\alpham}{\boldsymbol{\alpha}_m}
\newcommand{\betai}{\boldsymbol{\beta}^i}
\newcommand{\betavec}{\boldsymbol{\beta}}
\newcommand{\alphavec}{\boldsymbol{\alpha}}
\newcommand{\costvec}{\mathbf{C}}
\newcommand{\sigmoid}{\sigma}

\title{Explainer for the Logistic iSFS Model with Trust-Region Cost Updates}
\author{}
\date{\today}

\begin{document}

\maketitle
\tableofcontents

\section{Purpose of This Document}

This document explains the theoretical path used by the current implementation
of the incomplete-source feature-selection (iSFS) model in this repository.
The goal is to connect three things clearly:

\begin{enumerate}
  \item the original iSFS equations in Section 3 of \texttt{Multi\_source\_BWM.pdf},
  \item the logistic / binary-cross-entropy (BCE) version used in the code,
  \item the Python modules that implement each part of the method.
\end{enumerate}

The active modeling path in the repository is:

\begin{itemize}
  \item missingness profiles and overlapping training groups,
  \item logistic loss for binary classification,
  \item profile-specific source weights $\alpham$,
  \item shared source coefficients $\betai$,
  \item class-cost vector $\costvec$ updated by a trust-region step.
\end{itemize}

Two scope notes are important:

\begin{itemize}
  \item The original paper derives the optimization mainly with least squares.
  \item Remark 3 of the paper states that the iSFS model can be extended to
  logistic loss, and that is the path used here.
\end{itemize}

The trust-region cost block comes from \texttt{notes.tex}, but the loss inside
that cost block is adapted here from squared loss to BCE so that it matches the
current classifier.

\section{Where iSFS Starts in the Paper}

Equation (13) in the paper still belongs to the complete-data model from
Section 2. The incomplete-data iSFS model starts in Section 3, and the main
paper equations for the blockwise-missing setting are:

\begin{itemize}
  \item Equation (14): the iSFS optimization problem,
  \item Equation (15): the profile-level loss block,
  \item Equation (16): the $\alpha$ subproblem,
  \item Equation (17): the $\beta$ subproblem,
  \item Equation (18): the least-squares gradient for the $\beta$ block.
\end{itemize}

So, for the current blockwise-missing setting, Equations (14)--(18) are the
main source of truth.

\section{Original iSFS Formulation in the Paper}

The paper defines the incomplete-source feature-selection model as
\begin{equation}
\label{eq:paper14}
\min_{\alphavec, \betavec}
\frac{1}{|\pf|}
\sum_{m \in \pf}
f(\Xm, \betavec, \alpham, \ym)
 + \lambda R_{\beta}(\betavec)
\end{equation}
subject to
\begin{equation}
\label{eq:paper14-constraint}
R_{\alpha}(\alpham) \le 1,
\qquad \forall m \in \pf.
\end{equation}
This reproduces Equation (14) of the paper.

The paper then defines the profile-level loss block as
\begin{equation}
\label{eq:paper15}
f(\mathbf{X}, \betavec, \alphavec, \mathbf{y})
=
\frac{1}{n}
L\!\left(
\sum_{i=1}^{S} \alpha^i \mathbf{X}^i \betai,
\mathbf{y}
\right),
\end{equation}
which is Equation (15).

Here:
\begin{itemize}
  \item $S$ is the number of sources,
  \item $\pf$ is the set of profile-based training groups,
  \item $\Xm^i$ is the source-$i$ design matrix restricted to group $m$,
  \item $\betai$ is the shared coefficient vector for source $i$,
  \item $\alpha_m^i$ is the source weight for source $i$ inside group $m$,
  \item $L(\cdot,\cdot)$ is any convex loss,
  \item $R_{\alpha}$ and $R_{\beta}$ are regularization functions.
\end{itemize}

The profile-specific score is therefore
\begin{equation}
\label{eq:profile-score}
\zm
=
\sum_{i=1}^{S} \alpha_m^i \Xm^i \betai.
\end{equation}

This is the central linear model used throughout the iSFS formulation.

\section{From Least Squares to Logistic Loss}

In Section 3.2, the paper says that it focuses on the least-squares loss for
the derivation, but Remark 3 states that the same model can be extended to
logistic loss. That is the path used by the repository.

\subsection{Least-squares form in the paper}

For the least-squares case, the profile-level data-fit term is effectively
\begin{equation}
\label{eq:least-squares-profile}
\frac{1}{2n_m}
\left\|
\sum_{i=1}^{S} \alpha_m^i \Xm^i \betai
- \ym
\right\|_2^2.
\end{equation}

This is why the paper's $\beta$ block becomes quadratic.

\subsection{Logistic / BCE form used in this repository}

For binary classification, we instead use BCE written directly from logits:
\begin{equation}
\label{eq:bce-pointwise}
\ell(z, y)
=
\log(1 + e^z) - yz,
\qquad y \in \{0,1\}.
\end{equation}

For profile $m$, the logit vector is still
\begin{equation}
\label{eq:logit-vector-profile}
\zm
=
\sum_{i=1}^{S} \alpha_m^i \Xm^i \betai.
\end{equation}

So the logistic profile loss is
\begin{equation}
\label{eq:logistic-profile-loss}
f_m^{\mathrm{log}}(\alpham, \betavec)
=
\frac{1}{n_m}
\sum_{j \in G_m}
\ell(z_{m,j}, y_{m,j}).
\end{equation}

Two details matter here:
\begin{enumerate}
  \item the factor $\frac{1}{n_m}$ stays, because we still average over the
  samples in group $m$,
  \item the factor $\frac{1}{2}$ disappears, because it belonged to squared
  loss, not to BCE.
\end{enumerate}

This is the loss used in \texttt{loss\_functions.py}.

\section{Regularization Choice Used Here}

The paper leaves $R_{\alpha}$ and $R_{\beta}$ generic in Equation
\eqref{eq:paper14}, then states in Remark 4 and in the experiments section that
different regularizers are possible and that the $\ell_1$ penalty is used in
the experiments. The current repository follows that paper-aligned choice:
\begin{equation}
\label{eq:ralpha-choice}
R_{\alpha}(\alpham) = \|\alpham\|_1
\end{equation}
and
\begin{equation}
\label{eq:rbeta-choice}
R_{\beta}(\betavec)
=
\sum_{i=1}^{S} \|\betai\|_1.
\end{equation}

In code, this becomes:
\begin{itemize}
  \item an $\ell_1$-ball constraint for $\alpha$,
  \item an $\ell_1$ penalty plus proximal step for $\beta$.
\end{itemize}

This is implemented in \texttt{regularization.py}.

\section{Missingness Profiles and Block Construction}

This is the first major step in the incomplete-data model.

\subsection{Availability indicator for each participant}

Suppose there are $S$ sources. For participant $j$, define the binary source
availability vector
\begin{equation}
\label{eq:availability-vector}
\mathbf{I}_j = [I_{j1}, I_{j2}, \dots, I_{jS}],
\end{equation}
where
\begin{equation}
\label{eq:availability-entry}
I_{ji} =
\begin{cases}
1, & \text{if source } i \text{ is observed for participant } j, \\
0, & \text{if source } i \text{ is missing for participant } j.
\end{cases}
\end{equation}

The paper then converts this binary vector into a single integer profile:
\begin{equation}
\label{eq:profile-integer}
p_j
=
\sum_{i=1}^{S} I_{ji} 2^{S-i}.
\end{equation}

For $S=3$:
\begin{align*}
[1,1,1] &\mapsto 7, \\
[1,1,0] &\mapsto 6, \\
[1,0,1] &\mapsto 5, \\
[0,1,1] &\mapsto 3, \\
[1,0,0] &\mapsto 4, \\
[0,1,0] &\mapsto 2, \\
[0,0,1] &\mapsto 1.
\end{align*}

In code, we usually display the same idea as a binary string such as
\texttt{111}, \texttt{110}, \texttt{101}, and so on.

\subsection{Exact participant profile}

Each participant has exactly one exact observed-source profile:
\begin{equation}
\label{eq:exact-profile}
A_j = \{ i : I_{ji} = 1 \}.
\end{equation}

This exact profile is what prediction uses later.

\subsection{Overlapping optimization group $G_m$}

The paper does not train directly on exact profiles only. Instead, for each
source combination $m$, it builds an overlapping training group:
\begin{equation}
\label{eq:overlapping-group}
G_m
=
\{ j : (p_j \,\&\, m) = m \},
\end{equation}
where $\&$ denotes the bitwise AND operation.

This means:
\begin{itemize}
  \item sample $j$ belongs to $G_m$ if sample $j$ contains every source
  required by combination $m$,
  \item a sample has one exact profile, but it can belong to several training
  groups.
\end{itemize}

For example, if a participant has exact profile \texttt{111}, then that
participant belongs to all of the following training groups:
\[
\texttt{111}, \texttt{110}, \texttt{101}, \texttt{011}, \texttt{100},
\texttt{010}, \texttt{001}.
\]

This is why the training groups overlap.

\subsection{How the block matrices are formed}

Once $G_m$ is known, the block matrices are formed by row restriction:
\begin{equation}
\label{eq:block-matrix}
\Xm^i = \mathbf{X}^i[G_m, :],
\qquad
\ym = \mathbf{y}[G_m].
\end{equation}

So, for each combination $m$, we gather the rows whose available sources
contain $m$, and then use those rows to define the subproblem for that block.

\subsection{Small illustration table}

Assume there are 3 sources, denoted $S_1, S_2, S_3$.

\begin{longtable}{>{\raggedright\arraybackslash}p{0.08\textwidth}
                  >{\centering\arraybackslash}p{0.05\textwidth}
                  >{\centering\arraybackslash}p{0.05\textwidth}
                  >{\centering\arraybackslash}p{0.05\textwidth}
                  >{\raggedright\arraybackslash}p{0.14\textwidth}
                  >{\raggedright\arraybackslash}p{0.10\textwidth}
                  >{\raggedright\arraybackslash}p{0.10\textwidth}
                  >{\raggedright\arraybackslash}p{0.28\textwidth}}
\toprule
Sample & $S_1$ & $S_2$ & $S_3$ & Availability vector & Profile code & Exact profile & Training groups containing the sample \\
\midrule
\endhead
1 & 1 & 1 & 1 & $[1,1,1]$ & \texttt{111} & \texttt{111} & \texttt{001,010,011,100,101,110,111} \\
2 & 1 & 1 & 0 & $[1,1,0]$ & \texttt{110} & \texttt{110} & \texttt{100,010,110} \\
3 & 1 & 0 & 1 & $[1,0,1]$ & \texttt{101} & \texttt{101} & \texttt{100,001,101} \\
4 & 0 & 1 & 1 & $[0,1,1]$ & \texttt{011} & \texttt{011} & \texttt{010,001,011} \\
5 & 1 & 0 & 0 & $[1,0,0]$ & \texttt{100} & \texttt{100} & \texttt{100} \\
6 & 0 & 1 & 0 & $[0,1,0]$ & \texttt{010} & \texttt{010} & \texttt{010} \\
7 & 0 & 0 & 1 & $[0,0,1]$ & \texttt{001} & \texttt{001} & \texttt{001} \\
8 & 1 & 1 & 1 & $[1,1,1]$ & \texttt{111} & \texttt{111} & \texttt{001,010,011,100,101,110,111} \\
\bottomrule
\end{longtable}

Now take $m = 110$. Then
\begin{equation}
\label{eq:g110-example}
G_{110}
=
\{ j : (p_j \,\&\, 110) = 110 \}.
\end{equation}

In this small table, that includes samples whose exact profile is \texttt{110}
or \texttt{111}. So the rows in $G_{110}$ are the rows used to build
\[
\mathbf{X}_{110}^{(1)},
\qquad
\mathbf{X}_{110}^{(2)},
\qquad
\mathbf{y}_{110}.
\]

This is implemented in \texttt{missingness\_profiles.py}.

\section{Alpha Update When Beta and Cost Are Fixed}

The paper's Equation (16) says that when $\betavec$ is fixed, each
profile-specific weight vector $\alpham$ is obtained from
\begin{equation}
\label{eq:paper16}
\min_{\alphavec}
\left\|
\sum_{i=1}^{S} \alpha^i \mathbf{X}^i \betai
- \mathbf{y}
\right\|_2^2
\end{equation}
subject to
\begin{equation}
\label{eq:paper16-constraint}
R_{\alpha}(\alphavec) \le 1.
\end{equation}

\subsection{Logistic / cost-sensitive alpha objective used here}

In the current repository, with BCE and class costs, the profile-$m$ alpha
subproblem becomes
\begin{equation}
\label{eq:alpha-cost-objective}
f_m^C(\alpham)
=
\frac{1}{n_m}
\sum_{j \in G_m}
C_{y_{m,j}}
\ell(z_{m,j}, y_{m,j}),
\end{equation}
where
\begin{equation}
\label{eq:alpha-profile-logit}
z_{m,j}
=
\sum_{i=1}^{S}
\alpha_m^i
\big(\mathbf{x}_{m,j}^i\big)^T
\betai.
\end{equation}

\subsection{Gradient with respect to alpha}

Define the source-specific score vector
\begin{equation}
\label{eq:source-score}
\mathbf{s}_m^i = \Xm^i \betai.
\end{equation}

Then
\begin{equation}
\label{eq:z-as-source-scores}
\zm
=
\sum_{i=1}^{S} \alpha_m^i \mathbf{s}_m^i.
\end{equation}

The gradient of the BCE block with respect to the logits is
\begin{equation}
\label{eq:gm-logits}
\mathbf{g}_m
=
\frac{1}{n_m}
\left[
C_{\mathbf{y}_m}
\odot
\big(\sigmoid(\zm) - \ym\big)
\right].
\end{equation}

Therefore, by the chain rule,
\begin{equation}
\label{eq:alpha-gradient}
\frac{\partial f_m^C}{\partial \alpha_m^i}
=
\big(\mathbf{s}_m^i\big)^T \mathbf{g}_m
=
\big(\Xm^i \betai\big)^T \mathbf{g}_m.
\end{equation}

\subsection{Constrained update used in code}

Because the active choice is
\begin{equation}
\label{eq:alpha-l1-choice}
R_{\alpha}(\alpham) = \|\alpham\|_1,
\end{equation}
the code uses projected gradient:
\begin{equation}
\label{eq:alpha-temp}
\alpham^{\mathrm{temp}}
=
\alpham
- \eta_{\alpha} \nabla f_m^C(\alpham),
\end{equation}
\begin{equation}
\label{eq:alpha-projection}
\alpham^{\mathrm{next}}
=
\Pi_{\|\cdot\|_1 \le r_{\alpha}}
\big(
\alpham^{\mathrm{temp}}
\big).
\end{equation}

This is implemented in \texttt{alpha\_update.py}.

\section{Beta Update When Alpha and Cost Are Fixed}

The paper's Equation (17) says that when $\alphavec$ is fixed, the $\beta$
block is
\begin{equation}
\label{eq:paper17}
\min_{\betavec}
g(\betavec) + \lambda R_{\beta}(\betavec).
\end{equation}

For least squares, the paper writes
\begin{equation}
\label{eq:paper17-g}
g(\betavec)
=
\frac{1}{|\pf|}
\sum_{m \in \pf}
\frac{1}{2n_m}
\left\|
\sum_{i=1}^{S}
(\alpha_m^i \Xm^i)\betai
- \ym
\right\|_2^2.
\end{equation}

It then gives the least-squares gradient in Equation (18):
\begin{equation}
\label{eq:paper18}
\nabla g(\betai)
=
\frac{1}{|\pf|}
\sum_{m \in \pf}
\frac{1}{n_m}
I(m \,\&\, 2^{S-i} \ne 0)
(\alpha_m^i \Xm^i)^T
\left(
\sum_{r=1}^{S}
\alpha_m^r \Xm^r \boldsymbol{\beta}^r
- \ym
\right).
\end{equation}

\subsection{Logistic / cost-sensitive beta objective used here}

In the current repository, this becomes
\begin{equation}
\label{eq:beta-logistic-objective}
g_C^{\mathrm{log}}(\betavec)
=
\frac{1}{|\pf|}
\sum_{m \in \pf}
\frac{1}{n_m}
\sum_{j \in G_m}
C_{y_{m,j}}
\ell(z_{m,j}, y_{m,j}).
\end{equation}

\subsection{Logistic gradient with respect to beta}

The same logits gradient is
\begin{equation}
\label{eq:beta-gm}
\mathbf{g}_m
=
\frac{1}{n_m}
\left[
C_{\mathbf{y}_m}
\odot
\big(\sigmoid(\zm) - \ym\big)
\right].
\end{equation}

Then the gradient with respect to source-$i$ coefficients is
\begin{equation}
\label{eq:beta-gradient-logistic}
\nabla_{\betai} g_C^{\mathrm{log}}(\betavec)
=
\frac{1}{|\pf|}
\sum_{m : i \in m}
(\alpha_m^i \Xm^i)^T \mathbf{g}_m.
\end{equation}

This is the logistic analogue of Equation \eqref{eq:paper18}.

\subsection{Proximal update used in code}

Because
\begin{equation}
\label{eq:beta-l1-choice}
R_{\beta}(\betavec)
=
\sum_{i=1}^{S} \|\betai\|_1,
\end{equation}
the code uses a proximal-gradient step:
\begin{equation}
\label{eq:beta-v}
\mathbf{v}_i
=
\betai
- \eta_{\beta}
\nabla_{\betai} g_C^{\mathrm{log}}(\betavec),
\end{equation}
\begin{equation}
\label{eq:beta-prox}
\betai^{\mathrm{next}}
=
\operatorname{prox}_{\eta_{\beta}\lambda \|\cdot\|_1}(\mathbf{v}_i)
=
\operatorname{sign}(\mathbf{v}_i)
\odot
\max(|\mathbf{v}_i| - \eta_{\beta}\lambda, 0).
\end{equation}

This is implemented in \texttt{beta\_update.py}.

\section{Cost Vector and Trust-Region Update}

This is the main extension beyond the original paper path.

\subsection{Cost-sensitive BCE block}

For group $m$, the current repository uses
\begin{equation}
\label{eq:cost-sensitive-group-loss}
f_m^C(\Xm, \betavec, \alpham, \ym)
=
\frac{1}{n_m}
\sum_{j \in G_m}
C_{y_{m,j}}
\ell(z_{m,j}, y_{m,j}).
\end{equation}

So the full fixed-cost Block A objective is
\begin{equation}
\label{eq:block-a-objective}
\Phi(\alphavec, \betavec; \costvec)
=
\frac{1}{|\pf|}
\sum_{m \in \pf}
f_m^C(\Xm, \betavec, \alpham, \ym)
+
\lambda_{\beta}
\sum_{i=1}^{S} \|\betai\|_1
\end{equation}
subject to
\begin{equation}
\label{eq:block-a-constraint}
\|\alpham\|_1 \le r_{\alpha}.
\end{equation}

\subsection{Per-class loss vector}

With $\alphavec$ and $\betavec$ fixed, define
\begin{equation}
\label{eq:classwise-loss}
L_k(\alphavec, \betavec)
=
\frac{1}{|\pf|}
\sum_{m \in \pf}
\frac{1}{n_m}
\sum_{j \in G_m : y_{m,j}=k}
\ell(z_{m,j}, y_{m,j}),
\end{equation}
for each class $k$.

Then define the class-loss vector
\begin{equation}
\label{eq:class-loss-vector}
L(\alphavec, \betavec)
=
\big(
L_1(\alphavec, \betavec),
\dots,
L_K(\alphavec, \betavec)
\big)^T.
\end{equation}

The cost-only objective is then
\begin{equation}
\label{eq:phi-c}
\Phi_C(\costvec)
=
L(\alphavec, \betavec)^T \costvec.
\end{equation}

\subsection{Trust-region subproblem for the cost vector}

Let the current cost vector at outer iteration $t$ be $\costvec^{(t)}$, and let
$\costvec^0$ be the baseline cost vector. The trust-region step is written in
terms of $d_t$ where
\begin{equation}
\label{eq:trial-cost}
\costvec^{\mathrm{trial}} = \costvec^{(t)} + d_t.
\end{equation}

The trust-region problem solved for the $\costvec$ block is
\begin{equation}
\label{eq:trust-region-subproblem}
\max_{d}
\;
L_t^T d
\end{equation}
subject to
\begin{equation}
\label{eq:trust-region-constraints}
\frac{1}{2}\|d\|_2^2 \le \Delta_t,
\qquad
\costvec^{(t)} + d \ge 0,
\qquad
\frac{1}{2}\|\costvec^{(t)} + d - \costvec^0\|_2^2 \le \Delta_{\max}.
\end{equation}

This means:
\begin{itemize}
  \item the step must stay inside the current working trust region,
  \item the updated cost vector must stay nonnegative,
  \item the updated cost vector must remain inside the larger baseline-centered
  ball.
\end{itemize}

\subsection{Predicted increase, actual increase, and acceptance ratio}

Define the reduced objective
\begin{equation}
\label{eq:reduced-objective}
F(\costvec)
=
\min_{\alphavec, \betavec}
\Phi(\alphavec, \betavec; \costvec).
\end{equation}

Then the trust-region model predicts
\begin{equation}
\label{eq:predicted-increase}
\mathrm{Pred}_t = L_t^T d_t,
\end{equation}
while the actual change is
\begin{equation}
\label{eq:actual-increase}
\mathrm{Ared}_t
=
F(\costvec^{(t)} + d_t) - F(\costvec^{(t)}).
\end{equation}

The trust-region ratio is
\begin{equation}
\label{eq:trust-region-ratio}
\rho_t
=
\frac{\mathrm{Ared}_t}{\mathrm{Pred}_t}.
\end{equation}

If $\rho_t$ is large enough, the trial step is accepted. If not, the step is
rejected and the model keeps the previous $\costvec$, $\alphavec$, and
$\betavec$.

This is implemented in \texttt{costs.py} and orchestrated in
\texttt{train\_model.py}.

\section{Prediction Rule}

Training and prediction do not use the same grouping rule.

\subsection{During training}

Training uses the overlapping optimization groups $G_m$.

\subsection{During prediction}

Prediction uses the exact observed-source pattern of the new sample.

For a new sample $j$, define its exact observed-source set
\begin{equation}
\label{eq:prediction-exact-set}
A_j = \{ i : \text{source } i \text{ is observed for sample } j \}.
\end{equation}

Then prediction uses the exact-profile-specific alpha weights:
\begin{equation}
\label{eq:prediction-score}
z_j
=
\sum_{i \in A_j}
\alpha_{A_j}^i
\big(\mathbf{x}_j^i\big)^T
\betai.
\end{equation}

The probability of class 1 is
\begin{equation}
\label{eq:prediction-probability}
\Pr(y_j = 1 \mid x_j)
=
\sigmoid(z_j)
=
\frac{1}{1 + e^{-z_j}}.
\end{equation}

The predicted class is
\begin{equation}
\label{eq:prediction-threshold}
\hat{y}_j
=
\mathbf{1}\{\sigmoid(z_j) \ge \tau\},
\end{equation}
where $\tau = 0.5$ by default.

The learned cost vector $\costvec$ does not appear directly in the prediction
formula. Instead, $\costvec$ affects prediction indirectly by changing the
learned $\alphavec$ and $\betavec$ during training.

This is implemented in \texttt{prediction.py}.

\section{Full Training Loop Used in This Repository}

The current training loop is:

\subsection*{Step 0. Build profiles and training blocks}

From the source availability pattern:
\begin{enumerate}
  \item compute the exact profile of each sample,
  \item build the overlapping optimization groups $G_m$,
  \item build the block matrices $\Xm^i$ and labels $\ym$.
\end{enumerate}

\subsection*{Step 1. Initialize parameters}

The paper's Algorithm 2 says that each $\betai$ should be initialized by
fitting each source individually on the available data. The current code uses a
practical initialization:
\begin{itemize}
  \item initialize each $\betai$ with small random values,
  \item initialize each $\alpham$ with equal weight across the sources present
  in profile $m$,
  \item initialize the cost vector with
  \begin{equation}
  \label{eq:initial-cost}
  \costvec^{(0)} = [0.5, 0.5]^T
  \end{equation}
  and set the baseline cost vector to
  \begin{equation}
  \label{eq:baseline-cost}
  \costvec^0 = [0.5, 0.5]^T.
  \end{equation}
\end{itemize}

\subsection*{Step 2. Solve Block A with cost fixed}

With $\costvec^{(t)}$ fixed:
\begin{enumerate}
  \item solve the $\alpha$ subproblem with $\betavec$ fixed,
  \item solve the $\beta$ subproblem with $\alphavec$ fixed,
  \item repeat until the fixed-cost Block A objective stabilizes.
\end{enumerate}

Symbolically,
\begin{equation}
\label{eq:block-a-solve}
(\alphavec^{(t+1)}, \betavec^{(t+1)})
\approx
\arg\min_{\alphavec, \betavec}
\Phi(\alphavec, \betavec; \costvec^{(t)}).
\end{equation}

\subsection*{Step 3. Compute the class-loss vector}

With the updated $\alphavec^{(t+1)}$ and $\betavec^{(t+1)}$, compute
\begin{equation}
\label{eq:class-loss-step}
L^{(t+1)}
=
L(\alphavec^{(t+1)}, \betavec^{(t+1)}).
\end{equation}

\subsection*{Step 4. Solve the trust-region cost step}

Using $L^{(t+1)}$, solve for a trial cost step:
\begin{equation}
\label{eq:cost-step}
d_t
\approx
\arg\max_d L^{(t+1)T} d
\end{equation}
subject to the trust-region and nonnegativity constraints. This produces
\begin{equation}
\label{eq:trial-cost-step}
\costvec^{\mathrm{trial}} = \costvec^{(t)} + d_t.
\end{equation}

\subsection*{Step 5. Re-solve Block A at the trial cost}

Now fix $\costvec^{\mathrm{trial}}$ and solve again for the corresponding
$\alpha,\beta$ block:
\begin{equation}
\label{eq:block-a-trial}
(\alphavec^{\mathrm{trial}}, \betavec^{\mathrm{trial}})
\approx
\arg\min_{\alphavec, \betavec}
\Phi(\alphavec, \betavec; \costvec^{\mathrm{trial}}).
\end{equation}

\subsection*{Step 6. Accept or reject the cost step}

Compute
\begin{equation}
\label{eq:acceptance-terms}
\mathrm{Pred}_t = L_t^T d_t,
\qquad
\mathrm{Ared}_t = F(\costvec^{\mathrm{trial}}) - F(\costvec^{(t)}),
\qquad
\rho_t = \frac{\mathrm{Ared}_t}{\mathrm{Pred}_t}.
\end{equation}

Then:
\begin{itemize}
  \item if $\rho_t$ is large enough, accept the step and keep
  $(\costvec^{\mathrm{trial}}, \alphavec^{\mathrm{trial}}, \betavec^{\mathrm{trial}})$,
  \item otherwise reject the step and keep the previous
  $(\costvec^{(t)}, \alphavec^{(t+1)}, \betavec^{(t+1)})$.
\end{itemize}

\subsection*{Step 7. Update the trust-region radius}

Depending on $\rho_t$:
\begin{itemize}
  \item shrink the radius if the step was poor,
  \item keep the radius if the step was reasonable,
  \item enlarge the radius if the step was very reliable.
\end{itemize}

\subsection*{Step 8. Repeat until convergence}

Repeat Steps 2--7 until:
\begin{itemize}
  \item the outer BCD iteration limit is reached, or
  \item the cost step becomes too small, or
  \item the reduced objective stabilizes.
\end{itemize}

This full loop is implemented in \texttt{train\_model.py} and exposed by
\texttt{run\_training.py} and \texttt{run\_training.sh}.

\section{What Is Actually Minimized and What Is Maximized}

One subtle but important point is that the current training loop is not a pure
``minimize everything'' procedure.

For fixed cost $\costvec$, Block A solves
\begin{equation}
\label{eq:min-block-a}
\min_{\alphavec, \betavec}
\Phi(\alphavec, \betavec; \costvec).
\end{equation}

But for fixed $\alphavec$ and $\betavec$, the cost block solves
\begin{equation}
\label{eq:max-cost}
\max_{\costvec} \Phi_C(\costvec).
\end{equation}

So the outer scheme is a blockwise alternating procedure with:
\begin{itemize}
  \item minimization over $(\alphavec, \betavec)$,
  \item maximization over $\costvec$ inside the trust-region constraints.
\end{itemize}

This is why, during training output:
\begin{itemize}
  \item \texttt{Accepted=True} means the proposed trust-region move in
  $\costvec$ was kept,
  \item \texttt{Accepted=False} means the proposed move in $\costvec$ was
  rejected.
\end{itemize}

\section{Code Map}

The current code modules line up with the mathematics as follows.

\begin{center}
\begin{tabular}{ll}
\toprule
Theoretical part & Code file \\
\midrule
Stable logistic / BCE loss & \texttt{loss\_functions.py} \\
Profile construction and overlapping groups & \texttt{missingness\_profiles.py} \\
$\alpha$ subproblem & \texttt{alpha\_update.py} \\
$\beta$ subproblem & \texttt{beta\_update.py} \\
$\ell_1$ regularization operators & \texttt{regularization.py} \\
Cost-sensitive BCE and trust-region cost step & \texttt{costs.py} \\
Exact-profile prediction & \texttt{prediction.py} \\
Full outer BCD training loop & \texttt{train\_model.py} \\
User-facing runner and saved reports & \texttt{run\_training.py} \\
\bottomrule
\end{tabular}
\end{center}

\section{Short Summary}

The iSFS model in the paper is built around:
\begin{itemize}
  \item one shared coefficient vector $\betai$ per source,
  \item one profile-specific weight vector $\alpham$ per source combination,
  \item a profile-based decomposition of blockwise-missing data.
\end{itemize}

This repository keeps that structure, but replaces the least-squares data-fit
term with BCE from logits for binary classification, and adds a trust-region
update for the class-cost vector $\costvec$.

So the full current pipeline is:
\begin{enumerate}
  \item compute exact profiles and overlapping training groups,
  \item train the logistic iSFS model by alternating $\alpha$ and $\beta$,
  \item update the class-cost vector $\costvec$ by trust region,
  \item predict with the exact-profile logistic model.
\end{enumerate}

That is the theoretical path implemented by the current codebase.

\end{document}
